% GENERATED FILE. Edit canonical lessons, then run npm run generate:books.
% canonical-source-hash: fnv1a64:aaf1e913cc4178e5
% canonical-lessons: GU-C475-andh, GU-C475-baheru, GU-C475-mugu, GU-C475-najar, GU-C475-adakvu

\chapter{Seeing and Hearing}
\label{ch:gu-seeing-and-hearing}
\hlchaptermodality{475}

\begin{chapteropening}
\textbf{By the end of this chapter:} \emph{I can talk about sight, hearing and touch.}

Complete the last lesson of chapter 475: I can talk about sight, hearing and touch.
\end{chapteropening}

\section[\texorpdfstring{\gu{અંધ} (andh)}{અંધ (andh)}]{\gu{અંધ} (andh) — blind (the respectful word)}

\label{lesson:GU-C475-andh}



\begin{quote}
Before the new one: say the Gujarati for to spit, then the Gujarati for the palate.
\end{quote}

\subsection*{You'll want to know: \gu{અંધ}}
\textbf{\gu{અંધ}} — \emph{andh} — \textquotedblleft{}blind (the respectful word)\textquotedblright{}.

\begin{quote}
\gu{અરુણે} \gu{એક} \gu{અંધ} \gu{માણસને} \gu{રસ્તા} \gu{પર} \gu{મદદ} \gu{કરી}.
\end{quote}

\begin{quote}
\emph{Aruṇe ek andh māṇasne rastā par madad karī.}
\end{quote}

\begin{quote}
Arun helped a blind man in the street.
\end{quote}

\begin{cousinweb}
is a tatsama, Sanskrit \emph{andha}, \textquotedblleft{}blind\textquotedblright{}, as Monier-Williams records it; CDIAL lists the everyday inherited Gujarati forms \emph{ā̃dhũ} and \emph{ā̃dhḷũ} under the same word (CDIAL 385), but for a person the careful word is \textbf{\gu{અંધ}}.
\end{cousinweb}

\begin{grammarlens}[title={What you've built}]
One more word for describing what you lived and felt.
\end{grammarlens}

\subsection*{Your turn}
\begin{itemize}
\raggedright
  \item \emph{Say it:} \emph{andh}
  \item \emph{Say it:} \emph{andh}, once more
  \item \emph{Say it:} \emph{Aruṇe ek andh māṇasne rastā par madad karī.} — Arun helped a blind man in the street.
  \item \emph{Recall:} say the Gujarati for to spit, then the Gujarati for the palate, then say \emph{andh} again
\end{itemize}

\begin{tcolorbox}[breakable,colback=teal!4,colframe=teal!35!black,title={Before you move on}]
What does \gu{અંધ} mean? (\textquotedblleft{}Blind (the respectful word)\textquotedblright{}.) Say it once more.
\end{tcolorbox}
\section[\texorpdfstring{\gu{બહેરું} (baherũ)}{બહેરું (baherũ)}]{\gu{બહેરું} (baherũ) — deaf}

\label{lesson:GU-C475-baheru}



\begin{quote}
Before the new one: say the Gujarati for the palate, then the Gujarati for blind.
\end{quote}

\subsection*{You'll want to know: \gu{બહેરું}}
\textbf{\gu{બહેરું}} — \emph{baherũ} — \textquotedblleft{}deaf\textquotedblright{}.

\begin{quote}
\gu{દાદા} \gu{એક} \gu{કાને} \gu{બહેરા} \gu{છે}.
\end{quote}

\begin{quote}
\emph{dādā ek kāne baherā che.}
\end{quote}

\begin{quote}
Grandfather is deaf in one ear.
\end{quote}

\begin{cousinweb}
\textbf{\gu{બહેરું}} (\emph{baherũ}) is inherited from Sanskrit \emph{badhira}, \textquotedblleft{}deaf\textquotedblright{}; CDIAL lists the Gujarati word (CDIAL 9130).
\end{cousinweb}

\begin{grammarlens}[title={What you've built}]
One more word for describing what you lived and felt.
\end{grammarlens}

\subsection*{Your turn}
\begin{itemize}
\raggedright
  \item \emph{Say it:} \emph{baherũ}
  \item \emph{Say it:} \emph{baherũ}, once more
  \item \emph{Say it:} \emph{dādā ek kāne baherā che.} — Grandfather is deaf in one ear.
  \item \emph{Recall:} say the Gujarati for the palate, then the Gujarati for blind, then say \emph{baherũ} again
\end{itemize}

\begin{tcolorbox}[breakable,colback=teal!4,colframe=teal!35!black,title={Before you move on}]
What does \gu{બહેરું} mean? (\textquotedblleft{}Deaf\textquotedblright{}.) Say it once more.
\end{tcolorbox}
\section[\texorpdfstring{\gu{મૂંગું} (mū̃gũ)}{મૂંગું (mū̃gũ)}]{\gu{મૂંગું} (mū̃gũ) — silent; unable to speak}

\label{lesson:GU-C475-mugu}



\begin{quote}
Before the new one: say the Gujarati for blind, then the Gujarati for deaf.
\end{quote}

\subsection*{You'll want to know: \gu{મૂંગું}}
\textbf{\gu{મૂંગું}} — \emph{mū̃gũ} — \textquotedblleft{}silent; unable to speak\textquotedblright{}.

\begin{quote}
\gu{બીકથી} \gu{બાળક} \gu{મૂંગું} \gu{થઈ} \gu{ગયું}.
\end{quote}

\begin{quote}
\emph{bīkthī bāḷak mū̃gũ thaī gayũ.}
\end{quote}

\begin{quote}
The child went silent with fear.
\end{quote}

\begin{cousinweb}
\textbf{\gu{મૂંગું}} (\emph{mū̃gũ}) is inherited from Sanskrit \emph{mūka}, \textquotedblleft{}unable to speak\textquotedblright{}; CDIAL lists the Gujarati word (CDIAL 10231).
\end{cousinweb}

\begin{grammarlens}[title={What you've built}]
One more word for describing what you lived and felt.
\end{grammarlens}

\subsection*{Your turn}
\begin{itemize}
\raggedright
  \item \emph{Say it:} \emph{mū̃gũ}
  \item \emph{Say it:} \emph{mū̃gũ}, once more
  \item \emph{Say it:} \emph{bīkthī bāḷak mū̃gũ thaī gayũ.} — The child went silent with fear.
  \item \emph{Recall:} say the Gujarati for blind, then the Gujarati for deaf, then say \emph{mū̃gũ} again
\end{itemize}

\begin{tcolorbox}[breakable,colback=teal!4,colframe=teal!35!black,title={Before you move on}]
What does \gu{મૂંગું} mean? (\textquotedblleft{}Silent; unable to speak\textquotedblright{}.) Say it once more.
\end{tcolorbox}
\section[\texorpdfstring{\gu{નજર} (najar)}{નજર (najar)}]{\gu{નજર} (najar) — eyesight; a glance}

\label{lesson:GU-C475-najar}



\begin{quote}
Before the new one: say the Gujarati for deaf, then the Gujarati for silent.
\end{quote}

\subsection*{You'll want to know: \gu{નજર}}
\textbf{\gu{નજર}} — \emph{najar} — \textquotedblleft{}eyesight; a glance\textquotedblright{}.

\begin{quote}
\gu{ઉંમર} \gu{સાથે} \gu{દાદીની} \gu{નજર} \gu{નબળી} \gu{થઈ}.
\end{quote}

\begin{quote}
\emph{ummar sāthe dādīnī najar nabḷī thaī.}
\end{quote}

\begin{quote}
With age my grandmother's eyesight has grown weak.
\end{quote}

A feminine noun.

\begin{cousinweb}
\textbf{\gu{નજર}} (\emph{najar}) is usually traced to Arabic \emph{naẓar}, \textquotedblleft{}sight, a look\textquotedblright{}, which reached Gujarati through Persian.
\end{cousinweb}

\begin{grammarlens}[title={What you've built}]
One more word for describing what you lived and felt.
\end{grammarlens}

\subsection*{Your turn}
\begin{itemize}
\raggedright
  \item \emph{Say it:} \emph{najar}
  \item \emph{Say it:} \emph{najar}, once more
  \item \emph{Say it:} \emph{ummar sāthe dādīnī najar nabḷī thaī.} — With age my grandmother's eyesight has grown weak.
  \item \emph{Recall:} say the Gujarati for deaf, then the Gujarati for silent, then say \emph{najar} again
\end{itemize}

\begin{tcolorbox}[breakable,colback=teal!4,colframe=teal!35!black,title={Before you move on}]
What does \gu{નજર} mean? (\textquotedblleft{}Eyesight; a glance\textquotedblright{}.) Say it once more.
\end{tcolorbox}
\section[\texorpdfstring{\gu{અડકવું} (aḍakvũ)}{અડકવું (aḍakvũ)}]{\gu{અડકવું} (aḍakvũ) — to touch}

\label{lesson:GU-C475-adakvu}



\begin{quote}
Before the new one: say the Gujarati for silent, then the Gujarati for eyesight.
\end{quote}

\subsection*{You'll want to know: \gu{અડકવું}}
\textbf{\gu{અડકવું}} — \emph{aḍakvũ} — \textquotedblleft{}to touch\textquotedblright{}.

\begin{quote}
\gu{ગરમ} \gu{તપેલીને} \gu{અડકશો} \gu{નહીં}.
\end{quote}

\begin{quote}
\emph{garam tapelīne aḍakśo nahī̃.}
\end{quote}

\begin{quote}
Don't touch the hot pot.
\end{quote}

Past (he …): \gu{તે} \gu{ગરમ} \gu{તપેલીને} \gu{અડક્યો}, \emph{te garam tapelīne aḍakyo}.

\begin{cousinweb}
\textbf{\gu{અડકવું}} (\emph{aḍakvũ}) goes back to a reconstructed \emph{*aḍ-}, \textquotedblleft{}obstruct, stop\textquotedblright{}, which CDIAL connects with Dravidian; CDIAL lists the Gujarati \emph{aṛakvũ}, \textquotedblleft{}to touch\textquotedblright{} (CDIAL 187).
\end{cousinweb}

\begin{grammarlens}[title={What you've built}]
One more word for describing what you lived and felt.
\end{grammarlens}

\subsection*{Your turn}
\begin{itemize}
\raggedright
  \item \emph{Say it:} \emph{aḍakvũ}
  \item \emph{Say it:} \emph{aḍakvũ}, once more
  \item \emph{Say it:} \emph{garam tapelīne aḍakśo nahī̃.} — Don't touch the hot pot.
  \item \emph{Recall:} say the Gujarati for silent, then the Gujarati for eyesight, then say \emph{aḍakvũ} again
\end{itemize}

\begin{tcolorbox}[breakable,colback=teal!4,colframe=teal!35!black,title={Before you move on}]
What does \gu{અડકવું} mean? (\textquotedblleft{}To touch\textquotedblright{}.) Say it once more.
\end{tcolorbox}
